\documentclass[a4paper,12pt,twocolumn,landscape]{article}

\usepackage{FabZ}
\usepackage{vecteurs}
\usepackage{repere}

\usepackage{geometry}
\geometry{hmargin=0.5cm,vmargin=1.5cm}

\newcommand{\milieu}[6]{$\pointcoord{A_{\theenumi}}{#1}{#2}$ et $\pointcoord{B_{\theenumi}}{#3}{#4}$ \hfill \reponseEX{$\pointcoord{M_{\theenumi}}{#5}{#6}$}}

\newcommand{\extremite}[6]{$\pointcoord{A_{\theenumi}}{#1}{#2}$ et $\pointcoord{M_{\theenumi}}{#3}{#4}$ \hfill \reponseEX{$\pointcoord{B_{\theenumi}}{#5}{#6}$}}

\newcommand{\longueur}[5]{$\pointcoord{A_{\theenumi}}{#1}{#2}$ et $\pointcoord{B_{\theenumi}}{#3}{#4}$ \hfill \reponseEX{Longueur $A_{\theenumi}B_{\theenumi}$ = $#5$}}

\newcommand{\quadrilatere}[9]{$\pointcoord{A_{\theenumi}}{#1}{#2}$, $\pointcoord{B_\theenumi}{#3}{#4}$ $\pointcoord{C_{\theenumi}}{#5}{#6}$, $\pointcoord{D_\theenumi}{#7}{#8}$ \\ \textcolor{red}{$A_{\theenumi}B_{\theenumi}C_{\theenumi}D_{\theenumi}$ est un #9.}}

\newcommand{\parallelogrammeManqueA}[8]{$\pointcoord{B_{\theenumi}}{#1}{#2}$, $\pointcoord{C_{\theenumi}}{#3}{#4}$ $\pointcoord{D_{\theenumi}}{#5}{#6}$ \hfill \textcolor{red}{$\pointcoord{A_{\theenumi}}{#7}{#8}$}}

\newcommand{\parallelogrammeManqueB}[8]{$\pointcoord{A_{\theenumi}}{#1}{#2}$, $\pointcoord{C_{\theenumi}}{#3}{#4}$ $\pointcoord{D_{\theenumi}}{#5}{#6}$ \hfill \textcolor{red}{$\pointcoord{B_{\theenumi}}{#7}{#8}$}}

\newcommand{\parallelogrammeManqueC}[8]{$\pointcoord{A_{\theenumi}}{#1}{#2}$, $\pointcoord{B_{\theenumi}}{#3}{#4}$ $\pointcoord{D_{\theenumi}}{#5}{#6}$ \hfill \textcolor{red}{$\pointcoord{C_{\theenumi}}{#7}{#8}$}}

\newcommand{\parallelogrammeManqueD}[8]{$\pointcoord{A_{\theenumi}}{#1}{#2}$, $\pointcoord{B_{\theenumi}}{#3}{#4}$ $\pointcoord{C_{\theenumi}}{#5}{#6}$ \hfill \textcolor{red}{$\pointcoord{D_{\theenumi}}{#7}{#8}$}}

\newcommand{\rep}[3]{~$\left(#1,~#2,~#3\right)$}


%
%\setlength{\headheight}{0pt}
%\pagestyle{fancyplain}
%\fancyhf{}
%\lhead[]{\textbf{TD Vecteurs}}
%\chead[]{}
%\rhead[]{}
%
%\lfoot[]{}
%\cfoot[]{}
%\rfoot[]{}
%%\rfoot[]{Page \thepage~ sur \pageref{LastPage}}

\fancypagestyle{firststyle}
{
	\setlength{\headheight}{0em}
	\fancyhf{}
	\lhead[]{\textbf{Mathématiques\\DS1}}
	\chead[]{Seconde 6\\ vendredi 25 septembre 2015}
	\rhead[]{Séquence 1\\ Repérage dans le plan}

	\lfoot[]{}
	\cfoot[]{}
	\rfoot[]{}
%	\rfoot[]{Page \thepage~ sur \pageref{LastPage}}
}

\fancypagestyle{vide}
{
	\setlength{\headheight}{0em}
	\pagestyle{fancyplain}
	\def\headrulewidth{0em}
	\fancyhf{}
	\lhead[]{}
	\chead[]{}
	\rhead[]{}
	
	\lfoot[]{}
	\cfoot[]{}
	\rfoot[]{}
%	\rfoot[]{Page \thepage~ sur \pageref{LastPage}}
}

\usetikzlibrary{fadings}
\newcommand{\Fin}{node[xshift=-1.5ex,rotate=10]{F}
node[rotate=170]{i}
node[xshift=1.5ex,rotate=45]{n}}

\newcommand{\bonnesvacances}{node[xshift=0ex,rotate=10]{B}
node[xshift=1*1.5ex,rotate=170]{o}
node[xshift=2*1.5ex,rotate=45]{n}
node[xshift=3*1.5ex,rotate=128]{n}
node[xshift=4*1.5ex,rotate=75]{e}
node[xshift=5*1.5ex,rotate=130]{s}
node[xshift=6*1.5ex,rotate=85]{~}
node[xshift=7*1.5ex,rotate=128]{v}
node[xshift=8*1.5ex,rotate=43]{a}
node[xshift=9*1.5ex,rotate=4]{c}
node[xshift=10*1.5ex,rotate=145]{a}
node[xshift=11*1.5ex,rotate=5]{n}
node[xshift=12*1.5ex,rotate=25]{c}
node[xshift=13*1.5ex,rotate=105]{e}
node[xshift=14*1.5ex,rotate=45]{s}
}


\begin{document}

\vspace*{1em}

\begin{minipage}{0.45\textwidth}
\thispagestyle{firststyle}

\vspace*{1em}

\paragraph{Données pour les exercices 1 et 2 :}~\\
Le plan $\left(\mathscr{P}\right)$ est muni d'un repère~\emph{orthonormé}\rep{O}{I}{J}.\\ Dans ce repère, on considère les points suivants $\pointcoord{A}{-122}{-48}$ ; $\pointcoord{B}{28}{-98}$ et $\pointcoord{C}{-72}{102}$.
On désigne par $D$, $E$ et $F$ les milieux respectifs des segments $\left[AB\right]$,  $\left[BC\right]$ et $\left[AC\right]$.

\vspace*{2em}

\paragraph{Exercice~1}

\begin{enumerate}
	\item Déterminer la nature du triangle $ABC$.
	\item Déterminer les coordonnées des points $D$, $E$ et $F$ dans le repère\rep{O}{I}{J}.
	\item Montrer que le quadrilatère $ADEF$ est un carré.
	\item En déduire la nature du triangle $DEF$.
\end{enumerate}

\paragraph{Exercice~2} 

\begin{enumerate}
	\item Déterminer les coordonnées des points $A$, $B$, $C$, $D$, $E$ et $F$ dans le repère\rep{A}{B}{C}.
	\item Peut-on déterminer la longueur $AB$ dans ce repère~?\\Si oui, comment~? Sinon, pourquoi~?
	\item Déterminer les coordonnées des points $A$, $B$, $C$, $D$, $E$ et $F$ dans le repère\rep{A}{C}{B}.
\end{enumerate}

\vspace*{2em}

\hrule

\vspace*{2em}

\paragraph{Exercice~3} %Les questions de cet exercice sont indépendantes les unes des autres.
On se place dans le repère~\emph{orthonormé}\rep{O}{I}{J}. Dans ce repère, on considère les points $\pointcoord{A}{3}{0}$, $\pointcoord{B}{2}{-3}$ et $\pointcoord{C}{-10}{3}$.\\Vos réponses seront argumentées.

	\begin{enumerate}
		\item Le point $A$ appartient-il au cercle $\mathscr{C}$ de centre $C$ passant par~$B$~?
		\item Le point $C$ appartient-il à la médiatrice de $\left[AB\right]$~?
%		\item Déterminer les coordonnées de $Z$ pour que $AZBC$ soit un parallélogramme.
	\end{enumerate}



\vspace{-2em}


\end{minipage}
\newpage

\vspace*{1em}

\begin{minipage}{0.45\textwidth}
\thispagestyle{firststyle}

\vspace*{1em}

%\paragraph{Exercice~4} 
%
%On se place dans le repère \emph{orthonormé}\rep{O}{I}{J}. $ABCD$ est un rectangle de centre $O$. $A$ est tel que son abscisse est négative et son ordonnée positive. Le côté $AB$ est parallèle à l'axe des ordonnées. Pour chaque question, faire une figure en prenant $AB = 4 cm$ et $AD = 6 cm$. Donner les coordonnées des points $A$, $B$, $C$, $D$, $O$, $I$ et $J$~:
%	\begin{enumerate}
%		\item Dans le repère\rep{O}{I}{J}
%		\item Dans le repère\rep{B}{C}{A}
%		\item Dans le repère\rep{A}{B}{D}
%	\end{enumerate}

%\vspace{-2em}
\paragraph{Exercice~4} Un élève a voulu tracer un parallélogramme $ABCD$ mais il ne sait pas où placer le sommet $C$.\\ Comment pouvez-vous l'aider~?\\ Vous détaillerez votre démarche sur votre copie avant de réaliser des calculs.

\begin{center}

\begin{tikzpicture}[scale=1,every node/.style={scale=1}]

%\coordinate (X) at (-0.5,-0.5);
%\coordinate (Y) at (13,13);
%\clip (X) rectangle (Y);
\placerpoint{A}{1}{1}{below left};
\placerpoint{B}{10}{5}{above right};
\placerpoint{D}{4}{5}{above left};
%\placerpoint{M}{6}{4}{above left};
%\placerpoint{C}{11}{7}{below right};
%\draw[very thick,red] (A) -- (B) -- (C) -- (D) -- cycle;
\repereOIJ{-1}{11}{-1}{6};
\foreach \r in {5,10}
    \draw[thick, below right] (\r,0) node{\r};
\foreach \r in {5}
    \draw[thick, above left] (0,\r) node{\r};

\end{tikzpicture}

\end{center}

\vspace*{\stretch{1}}
\centering
\begin{tikzpicture}[scale=2,transform shape]
\draw (0,0) \Fin;
\draw (-1em,-1ex) -- (1em,-1ex);
\path[scope fading=south] (-1em,-0.25em) rectangle (1em,-3.75ex);
\draw[yscale=-1] (0,2ex) \Fin;
\end{tikzpicture}
\vspace*{\stretch{1}}


\end{minipage}
\end{document}
%
%
%\paragraph{Exercice~2} Propriétés des quadrilatères
%
%\begin{tikzpicture}[scale=1,every node/.style={scale=0.7}]
%\tikzstyle{debutfin}=[ellipse,draw,text=red]
%\tikzstyle{instruct}=[rectangle,draw,fill=yellow!50]
%\tikzstyle{test}=[diamond, aspect=6,thick,
%draw=blue,fill=yellow!50,text=blue]
%\tikzstyle{es}=[rectangle,draw,rounded corners=4pt,fill=blue!25]
%
%\node[debutfin] (debut) at (0,3) {Début};
%\node[es] (lire) at (0,2) {Prendre un quadrilatère $Q$};
%\node[test] (test) at (0,0) {Les diagonales ont-elles un milieu commun \ ?};
%%\node[instruct] (init) at (-2,2.5) {$S\leftarrow 0$};
%\node[instruct] (plus) at (0,-2.5) {$S\leftarrow S+N$};
%\node[instruct] (moins) at (0,-3.5) {$N\leftarrow N-1$};
%\node[es] (afficher) at (-4,-2) {Afficher la somme $S$};
%\node[debutfin] (fin) at (-4,-3) {Fin};
%
%\tikzstyle{suite}=[->,>=stealth,thick,rounded corners=4pt]
%\draw[suite](debut) -- (lire);
%\draw[suite](lire) -- (test.north);
%%\draw[suite](init) -- (test.north);
%\draw[suite](test.south) -- (plus);
%\draw[suite](plus) -- (moins);
%\draw[suite](test) -| (afficher);
%\draw[suite](afficher) -- (fin);
%\end{tikzpicture}
%
%
%%%%%\begin{tikzpicture}
%%%%%% définition des styles
%%%%%\tikzstyle{quadri}=[rectangle,draw,fill=yellow!50,text=blue]
%%%%%\tikzstyle{estun}=[->,>=latex,very thick,dotted]
%%%%%% les nœuds
%%%%%\node[quadri] (Q) at (0,3) {Quadrilatère};
%%%%%\node[quadri] (P) at (0,1.5) {Parallélogramme};
%%%%%\node[quadri] (R) at (-3,0) {Rectangle};
%%%%%\node[quadri] (L) at (3,0) {Losange};
%%%%%\node[quadri] (C) at (5,-1.5) {Carré};
%%%%%% les flèches
%%%%%\draw[estun] (P)--(Q);
%%%%%\draw[estun] (R)to[bend left](Q); \draw[estun] (R)--(P);
%%%%%\draw[estun] (L)--(Q.south east); \draw[estun] (L)--(P);
%%%%%\draw[estun] (C)to[bend right](Q.east); \draw[estun] (C)to[bend left](P);
%%%%%\draw[estun] (C)--(L.south east); \draw[estun] (C)to[bend left](R);
%%%%%% la légende
%%%%%\draw[estun] (-4.5,2.5)--(-3,2.5)node[midway,above]{est un};
%%%%%\end{tikzpicture}
